% Source: https://artofproblemsolving.com/community/c3h1161680_synthetic_division
% Stable identity: aops-topic-1161680-q01
Factor the divisor: $x^2-x-2=(x-2)(x+1)$.

\textbf{Divide by $x-2$.} Synthetic division with $2$ on the coefficients $1,-3,2,5,-1$ gives $1,\,-1,\,0,\,5$ and remainder $9$:
\[
P(x)=(x-2)(x^3-x^2+5)+9 .
\]
\textbf{Divide the quotient by $x+1$.} Synthetic division with $-1$ on $1,-1,0,5$ gives $1,\,-2,\,2$ and remainder $3$:
\[
x^3-x^2+5=(x+1)(x^2-2x+2)+3 .
\]
\textbf{Combine.}
\[
P(x)=(x-2)\bigl[(x+1)(x^2-2x+2)+3\bigr]+9=(x^2-x-2)(x^2-2x+2)+3(x-2)+9 .
\]
So the quotient is $\boxed{x^2-2x+2}$ and the remainder is $\boxed{3x+3}$.

Check with the roots of the divisor: $P(2)=16-24+8+10-1=9=3\cdot2+3$ and $P(-1)=1+3+2-5-1=0=3(-1)+3$. In general, dividing by $(x-a)(x-b)$ this way gives quotient $Q_2$ and remainder $r_2(x-a)+r_1$, where $r_1$, $r_2$ are the two synthetic-division remainders.
